% !TeX root = ../navier-stokes-manuscript.tex
% ============================================================
% 2.1 Vortex Stretching
% ============================================================

\subsection{Vortex Stretching}\label{sub:ThePerfectlyStretchingVortex}

One narrow material vortex tube, stretched along its axis, lengthens and squeezes its cross-section. Rotating fluid moves inward as the tube narrows, and the interior spins faster. This spin-up is the positive stretching contribution. Whether the rotation strengthens overall still depends on its complete signed balance with viscosity inside the same velocity evolution:
\[
  \underbrace{\partial_t u}_{\text{local acceleration}}
  +
  \underbrace{(u\cdot\nabla)u}_{\text{nonlinear transport}}
  =
  \underbrace{-\nabla p}_{\text{pressure}}
  +
  \underbrace{\nu\Delta u}_{\text{viscous diffusion}}.
\]
The tube remains local, and its net rotational change remains unresolved. Suppose its strain is locally uniform and axisymmetric across the tube, with vorticity aligned along its axis. If \(e\) names that axis and \(L(t)\) the tube's length, the axial strain fixes the rate of extension:
\[
  S:=\frac12\bigl(\nabla u+(\nabla u)^{\mathsf T}\bigr),
  \qquad
  Se=\sigma(t)e,
  \qquad
  \sigma(t)>0,
  \qquad
  \frac{\dot L(t)}{L(t)}
  =
  e\cdot Se
  =
  \sigma(t).
\]
Positive \(\sigma(t)\) lengthens the tube. Incompressibility forces its transverse area to contract at the same time. Its material volume \(\mathcal V\) stays fixed, so define its effective transverse area by \(A(t):=\mathcal V/L(t)\) and its effective radius by \(r_{\mathrm{eff}}(t):=\sqrt{A(t)/\pi}\). Their rates are
\[
  \nabla\cdot u=0,
  \qquad
  \frac{d}{dt}(A L)=0,
  \qquad
  \frac{\dot A}{A}=-\sigma(t),
  \qquad
  \frac{\dot r_{\mathrm{eff}}}{r_{\mathrm{eff}}}
  =
  -\frac{\sigma(t)}{2}.
\]
The effective radius now contracts at half the axial logarithmic rate. The vorticity \(\omega:=\nabla\times u\), carried by this same motion, evolves by
\[
  \frac{D\omega}{Dt}
  =
  S\omega+\nu\Delta\omega,
  \qquad
  \frac{D}{Dt}
  :=
  \partial_t+u\cdot\nabla,
\]
and the alignment \(\omega=|\omega|e\) turns the stretching part into a positive multiple of the vorticity:
\[
  S\omega
  =
  \sigma(t)\omega,
  \qquad
  \frac12\frac{D|\omega|^2}{Dt}
  =
  \sigma(t)|\omega|^2
  +
  \nu\,\omega\cdot\Delta\omega.
\]
The first term pushes the rotation upward; the viscous term can oppose it. On any interval where their signed sum is positive, rotation strengthens while the radius contracts. That local concentration must still be reconciled with the global kinetic energy. The identity proved in \Cref{prop:ClassicalKineticEnergyIdentity}, together with the antisymmetric part of \(\nabla u\), gives
\[
  \norm{u(t)}_2
  \leq
  \norm{u_0}_2
  <\infty,
  \qquad
  \left|\frac12\bigl(\nabla u-(\nabla u)^{\mathsf T}\bigr)\right|_{\mathrm F}
  =
  \frac{|\omega|}{\sqrt2}
  \leq
  |\nabla u|_{\mathrm F}.
\]
Finite kinetic energy does not close off a rising local rotational part of the gradient. The finer-scale question needs a complete solution for comparison.

\divider{\NavierStokes{} Scaling}

At a time \(t_0<T_*\), call the tube's effective radius \(\ell:=r_{\mathrm{eff}}(t_0)\). For \(\lambda>1\), define the rescaled velocity and pressure by
\[
  u_\lambda(x,t)
  :=
  \lambda u(\lambda x,\lambda^2t),
  \qquad
  p_\lambda(x,t)
  :=
  \lambda^2p(\lambda x,\lambda^2t).
\]
This pair is a complete \NavierStokes{} solution at another scale. At time \(\lambda^{-2}t_0\), its corresponding segment has effective radius \(\lambda^{-1}\ell\); it is not the opening tube at a later material time. The relative strength of the four mechanisms is decided by how each one transforms:
\[
\begin{aligned}
  \partial_tu_\lambda(x,t)
  &=
  \lambda^3(\partial_tu)(\lambda x,\lambda^2t),
  \\
  (u_\lambda\cdot\nabla)u_\lambda(x,t)
  &=
  \lambda^3\bigl((u\cdot\nabla)u\bigr)(\lambda x,\lambda^2t),
  \\
  \nabla p_\lambda(x,t)
  &=
  \lambda^3(\nabla p)(\lambda x,\lambda^2t),
  \\
  \nu\Delta u_\lambda(x,t)
  &=
  \lambda^3\nu(\Delta u)(\lambda x,\lambda^2t).
\end{aligned}
\]
Every momentum term gains the same cubic factor. Finer scale gives viscosity no relative advantage over nonlinear transport, while pressure and local acceleration remain in the same balance.

\divider{Critical Height}

The equation keeps its balance, but the size of the rescaled velocity depends on which homogeneous Sobolev level measures it. Let \(\Lambda:=(-\Delta)^{1/2}\). Its Fourier transform is
\[
  \widehat{u_\lambda}(\xi,t)
  =
  \lambda^{-2}\widehat u(\lambda^{-1}\xi,\lambda^2t),
\]
and the corresponding change of variables produces the exponent that separates those levels:
\[
  \norm{u_\lambda(t)}_{\dot H^s}^2
  =
  \lambda^{2s-1}
  \norm{u(\lambda^2t)}_{\dot H^s}^2.
\]
At \(s=0\), kinetic energy loses one power of \(\lambda\). At \(s=1\), the first derivatives gain one. The exponent \(2s-1\) vanishes at \(s=\tfrac12\), so define the scale-neutral height by
\[
  H(t)
  :=
  \frac12\norm{u(t)}_{\dot H^{1/2}}^2
  =
  \frac12\norm{\Lambda^{1/2}u(t)}_2^2.
\]
Writing \(H_\lambda\) for this quantity evaluated on \(u_\lambda\), the three levels separate explicitly:
\[
  \norm{u_\lambda(t)}_2^2
  =
  \lambda^{-1}\norm{u(\lambda^2t)}_2^2,
  \qquad
  H_\lambda(t)=H(\lambda^2t),
  \qquad
  \norm{\nabla u_\lambda(t)}_2^2
  =
  \lambda\norm{\nabla u(\lambda^2t)}_2^2.
\]
The rescaled solution carries less kinetic energy and a larger first-derivative norm, while \(H\) stays unchanged by the spatial rescaling. Its time derivative still has no fixed sign.

\divider{Nonlinear Production versus Viscous Dissipation}

At critical height, give the nonlinear contribution its signed value
\[
  \mathcal P_{1/2}[u](t)
  :=
  -\operatorname{Re}
  \pair
    {\Lambda^{1/2}u(t)}
    {\Lambda^{1/2}\bigl((u(t)\cdot\nabla)u(t)\bigr)}_{L^2}.
\]
The pressure pairing vanishes because \(\nabla\cdot u=0\), while the Laplacian contributes \(-\nu\norm{\Lambda^{3/2}u}_2^2\). The unresolved derivative now has the complete balance
\[
  H'(t)
  +
  \nu\norm{\Lambda^{3/2}u(t)}_2^2
  =
  \mathcal P_{1/2}[u](t).
\]
Thus \(H\) rises exactly when signed nonlinear production exceeds positive viscous dissipation. Rescaling changes both quantities by
\[
  \mathcal P_{1/2}[u_\lambda](t)
  =
  \lambda^2\mathcal P_{1/2}[u](\lambda^2t),
  \qquad
  \nu\norm{\Lambda^{3/2}u_\lambda(t)}_2^2
  =
  \lambda^2\nu\norm{\Lambda^{3/2}u(\lambda^2t)}_2^2.
\]
Their common quadratic factor preserves their relative strength. It also leaves open the physical time available for viscosity to act between two widths.

\divider{The Heat Clock}

A scale-localized velocity structure of width \(r\) is concentrated near frequencies \(|\xi|\simeq r^{-1}\). Under the comparison equation \(v_t=\nu\Delta v\), its Fourier modes change by the heat multiplier
\[
  \widehat v(\xi,t+\delta t)
  =
  e^{-\nu|\xi|^2\delta t}\widehat v(\xi,t).
\]
The multiplier changes by order one once its exponent reaches order one:
\[
  \nu|\xi|^2\delta t\simeq1,
\]
so the linear viscous comparison time for width \(r\) is
\[
  \tau_\nu(r)
  :=
  \frac{r^2}{\nu}.
\]
This clock measures diffusion at that localization width. Halving the width quarters the time, and an arbitrary refinement gives
\[
  \tau_\nu(\lambda^{-1}r)
  =
  \lambda^{-2}\tau_\nu(r).
\]
The same inverse-square contraction appears in the full \NavierStokes{} rescaling. Repeating it dyadically asks whether the shrinking heat times can fit inside finite time.

\divider{The Zeno Cascade}

For the localization widths and their heat clocks, set
\[
  r_n:=2^{-n}r_0,
  \qquad
  \tau_n:=\frac{r_n^2}{\nu},
\]
which gives the finite geometric sum
\[
  \tau_n
  =
  \frac{r_0^2}{\nu}4^{-n},
  \qquad
  \sum_{n=0}^{\infty}\tau_n
  =
  \frac{4r_0^2}{3\nu}
  <
  \infty.
\]
The arithmetic becomes a possible fluid history only if one \NavierStokes{} velocity field exhibits scale-localized structures at widths
\[
  r_0>r_1>r_2>\cdots\longrightarrow0
\]
and at sampled times
\[
  t_0<t_1<t_2<\cdots\uparrow T_*<\infty,
  \qquad
  t_{n+1}-t_n\asymp\tau_n,
\]
with comparison constants independent of \(n\). Under that one-field condition, its remaining time is
\[
  T_*-t_n
  \asymp
  \sum_{k=n}^{\infty}\tau_k
  =
  \frac{4r_n^2}{3\nu}.
\]
At the \(n\)-th sampled width, both the next gap and the total time left are comparable to \(r_n^2/\nu\). The spatial localization changes on intervals that collapse to zero as \(t\uparrow T_*\). Those intervals do not decide what happens to this same field's first and successive Eulerian time derivatives at one fixed spatial point. Their behavior is the question left open by the cascade.
